% =============================================================
%  Module Description — Data Structures and Algorithms II (CS232-F27)
%  DRAFT for a planned module (Semester 3 of the Foundation programme).
%  The code CS232-F27 is a proposal (semester digit + position in the
%  semester); add the term suffix (e.g. -S27) once the term is known and
%  make it equal the course's `code` in src/content/courses/.
%  Compile with: pdflatex cs232-module-description.tex
%  Write "TODO(owner): ..." where information is still missing;
%  the website highlights it.
% =============================================================
\documentclass[11pt,a4paper]{article}

\usepackage[utf8]{inputenc}
\usepackage[T1]{fontenc}
\usepackage{lmodern}
\usepackage[margin=2.2cm]{geometry}
\usepackage{booktabs}
\usepackage{tabularx}
\usepackage{array}
\usepackage{enumitem}
\usepackage{xcolor}
\usepackage{titlesec}
\usepackage{fancyhdr}
\usepackage[hidelinks]{hyperref}

% ---------- Module details (edit these) ----------------------
% \ModuleCode must equal the course's `code` in src/content/courses/.
\newcommand{\ModuleTitle}{Data Structures and Algorithms II}
\newcommand{\ModuleCode}{CS232-F27}
\newcommand{\Programme}{Foundation programme}
\newcommand{\Level}{Semester 3}
\newcommand{\Semester}{Fall 2027}
\newcommand{\Credits}{7 ECTS}
\newcommand{\ModuleLeader}{Dr. Ali Khudiyev}
\newcommand{\Email}{ali.khudiyev@wosm.academy}
\newcommand{\OfficeHours}{Monday--Friday, 10:00--17:00 (Baku time, UTC+4)}
\newcommand{\Prerequisites}{CS103-S27 Data Structures and Algorithms I}

% ---------- Styling ------------------------------------------
\definecolor{accent}{RGB}{0,70,130}
\titleformat{\section}{\large\bfseries\color{accent}}{\thesection.}{0.5em}{}
\titlespacing*{\section}{0pt}{1.2em}{0.5em}
\setlist{itemsep=2pt, topsep=4pt}
\setlength{\parindent}{0pt}
\setlength{\parskip}{4pt}

\pagestyle{fancy}
\fancyhf{}
\lhead{\small \ModuleCode\ -- \ModuleTitle}
\rhead{\small Module Description}
\cfoot{\small \thepage}
\renewcommand{\headrulewidth}{0.4pt}

% =============================================================
\begin{document}

\begin{center}
  {\LARGE\bfseries\color{accent} \ModuleTitle}\\[4pt]
  {\large Module Description}
\end{center}
\vspace{6pt}

% ---------- Summary table (not shown on the website: the course page
% shows these details from the macros above) --------------------------
\begin{tabularx}{\textwidth}{>{\bfseries}l X >{\bfseries}l X}
\toprule
Module code   & \ModuleCode    & Credits      & \Credits \\
Programme     & \Programme     & Level        & \Level \\
Semester      & \Semester      & Prerequisites & \Prerequisites \\
Module leader & \ModuleLeader  & Email        & \href{mailto:\Email}{\Email} \\
Office hours  & \multicolumn{3}{l}{\OfficeHours} \\
\bottomrule
\end{tabularx}

% ---------- Sections (shown on the website) ------------------
\section{Module Overview}
Data Structures and Algorithms II continues from Data Structures and Algorithms I.
Where the first module concentrated on data structures and sorting, this module
concentrates on \emph{algorithm design}: graph algorithms, and the general techniques
of divide and conquer, greedy algorithms, dynamic programming and network flow.

The last part of the module looks at problems for which no efficient algorithm is
known. You will learn to recognise NP-hard problems through reductions and to cope
with them using approximation, randomisation and systematic search. This prepares
you for Theory of Computation in Semester 4. Programming is done in C, as in Data
Structures and Algorithms I.

\section{Aims}
\begin{itemize}
  \item Teach the main paradigms of algorithm design and how to justify algorithms built with them.
  \item Give you the standard algorithms on graphs and networks and practice in modelling problems as graphs.
  \item Introduce amortised analysis and advanced data structures used inside efficient algorithms.
  \item Introduce intractability and practical strategies for hard problems.
\end{itemize}

\section{Learning Outcomes}
After completing this course, you will be able to:
\begin{enumerate}[label=LO\arabic*.]
  \item Model problems as graphs and apply and analyse algorithms for traversal, topological sorting, connectivity, shortest paths and minimum spanning trees.
  \item Design algorithms with divide and conquer, greedy choice and dynamic programming, and justify them with recurrences, exchange arguments and optimal substructure.
  \item Formulate problems as network flow and bipartite matching problems and solve them.
  \item Analyse algorithms and data structures, such as union--find, with aggregate, accounting and potential-method amortised analysis.
  \item Prove problems NP-complete by polynomial-time reduction and choose a suitable strategy for them: approximation, randomisation, or exact search.
\end{enumerate}

\section{Syllabus}
\begin{tabularx}{\textwidth}{c X}
\toprule
\textbf{Unit} & \textbf{Topic} \\
\midrule
1  & Graphs: representations, breadth-first and depth-first search \\
2  & Applications of DFS: topological sorting, strongly connected components \\
3  & Shortest paths I: Dijkstra's algorithm \\
4  & Shortest paths II: Bellman--Ford; all-pairs shortest paths \\
5  & Minimum spanning trees: Prim and Kruskal; union--find \\
6  & Amortised analysis: aggregate, accounting and potential methods \\
7  & Divide and conquer: Karatsuba, Strassen, closest pair of points \\
8  & Greedy algorithms: exchange arguments, scheduling, Huffman coding \\
9  & Dynamic programming I: principles; sequence problems \\
10 & Dynamic programming II: knapsack, edit distance, shortest paths revisited \\
11 & Network flow: Ford--Fulkerson, Edmonds--Karp, max-flow min-cut theorem \\
12 & Applications of flow: bipartite matching and related problems \\
13 & String algorithms: Knuth--Morris--Pratt, Rabin--Karp, suffix arrays \\
14 & Reductions, P, NP and NP-completeness \\
15 & Approximation algorithms \\
16 & Randomised algorithms; backtracking and branch and bound \\
17 & Course Review \& Exam \\
\bottomrule
\end{tabularx}

\section{Teaching and Learning Methods}
The module is self-paced. All materials are provided on Moodle.
Each unit has a lecture, exercises on designing and proving algorithms, and
programming tasks in C.

\begin{tabularx}{0.6\textwidth}{X r}
\toprule
\textbf{Activity} & \textbf{Hours} \\
\midrule
Lectures            & 34 \\
Labs / Tutorials    & 36 \\
Independent study   & 130 \\
\midrule
\textbf{Total}      & \textbf{200} \\
\bottomrule
\end{tabularx}

\section{Assessment}
Your knowledge will be assessed throughout the course with quizzes, programming
assignments and the final exam. The exam checks whether you can choose a suitable
design technique for an unfamiliar problem, develop an algorithm with it, and argue
its correctness and running time.

\begin{tabularx}{\textwidth}{X c c c}
\toprule
\textbf{Component} & \textbf{Weight} & \textbf{Due} & \textbf{LOs} \\
\midrule
Quizzes & 20\% & Throughout & LO1--LO5 \\
Programming assignments & 40\% & Throughout & LO1--LO3 \\
Final exam & 40\% & End of module & LO1--LO5 \\
\midrule
\textbf{Total} & \textbf{100\%} & & \\
\bottomrule
\end{tabularx}

\section{Reading List}
\textbf{Essential}
\begin{itemize}
  \item T.\ H.\ Cormen, C.\ E.\ Leiserson, R.\ L.\ Rivest, C.\ Stein: \textit{Introduction to Algorithms}, 4th ed., MIT Press, 2022.
  \item J.\ Kleinberg, \'E.\ Tardos: \textit{Algorithm Design}, Addison-Wesley, 2005.
\end{itemize}
\textbf{Recommended}
\begin{itemize}
  \item S.\ Dasgupta, C.\ H.\ Papadimitriou, U.\ V.\ Vazirani: \textit{Algorithms}, McGraw-Hill, 2006.
  \item S.\ S.\ Skiena: \textit{The Algorithm Design Manual}, 3rd ed., Springer, 2020.
\end{itemize}

\section{Materials and Licence}
Course materials are published under the
\href{https://creativecommons.org/licenses/by-sa/4.0/}{Creative Commons Attribution-ShareAlike 4.0}
licence unless stated otherwise. See \href{https://wosm.academy/licence/}{wosm.academy/licence} for details.

\section{Policies}
\begin{itemize}
  \item \textbf{Attendance:} Attendance is not mandatory. The module is self-paced, so you can study the materials when it suits you.
  \item \textbf{Late submissions:} Work submitted after the deadline is penalised.
  \item \textbf{Academic integrity:} All submitted work must be your own.
        See the \href{https://wosm.academy/code-of-conduct/}{Code of Conduct}.
\end{itemize}

\end{document}
